% !TEX program = xelatex
% ============================================================
%  خطة إصلاح مطابقة مصطلحات القاموس مع النص
%  (حدود الكلمات · التسامح الصرفي · توحيد المنطق · صحة التعريف)
%  الاستخدام على Overleaf:
%    1) New Project ← Blank Project
%    2) الصق المحتوى في main.tex
%    3) Menu ← Compiler ← اختر: XeLaTeX   (مهم للعربية)
%    4) Recompile ← نزّل PDF
% ============================================================
\documentclass[11pt]{article}

\usepackage{polyglossia}
\setmainlanguage{arabic}
\setotherlanguage{english}
\newfontfamily\arabicfont[Script=Arabic,Scale=1.0]{Amiri}
\newfontfamily\arabicfontsf[Script=Arabic,Scale=1.0]{Amiri}
\setmonofont{DejaVu Sans Mono}[Scale=0.8]

\usepackage{geometry}
\geometry{a4paper, margin=2.2cm}

\usepackage{xcolor}
\usepackage{graphicx}
\usepackage{booktabs}
\usepackage{longtable}
\usepackage{array}
\usepackage{enumitem}
\usepackage{titlesec}
\usepackage{fancyhdr}
\usepackage{hyperref}
\usepackage{tcolorbox}
\tcbuselibrary{skins,breakable}
\usepackage{fontawesome5}

% ── الألوان (هوية يَسِّر) ──
\definecolor{gold}{HTML}{B8933B}
\definecolor{teal}{HTML}{0F3D3E}
\definecolor{darkbg}{HTML}{101414}
\definecolor{okgreen}{HTML}{27AE60}
\definecolor{warnred}{HTML}{C0392B}
\definecolor{softgray}{HTML}{F4F4F4}
\definecolor{midgray}{HTML}{7A7A7A}
\definecolor{codebg}{HTML}{1A2020}

\pagestyle{fancy}
\fancyhf{}
\renewcommand{\headrulewidth}{0.4pt}
\renewcommand{\footrulewidth}{0.4pt}
\fancyhead[R]{\small\color{gold}خطة مطابقة مصطلحات القاموس — منصة يَسِّر}
\fancyhead[L]{\small\color{midgray}Glossary Matching Fix}
\fancyfoot[C]{\small\thepage}

\titleformat{\section}{\Large\bfseries\color{gold}}{\thesection}{0.6em}{}
\titleformat{\subsection}{\large\bfseries\color{darkbg}}{\thesubsection}{0.5em}{}
\titleformat{\subsubsection}{\normalsize\bfseries\color{midgray}}{\thesubsubsection}{0.4em}{}

\setlist[itemize]{rightmargin=0pt, itemsep=2pt, topsep=3pt}
\setlist[enumerate]{rightmargin=0pt, itemsep=2pt, topsep=3pt}

\hypersetup{
  colorlinks=true, linkcolor=gold, urlcolor=gold,
  pdftitle={خطة مطابقة مصطلحات القاموس}, pdfauthor={منصة يَسِّر}
}

% ── صناديق مخصّصة ──
\newtcolorbox{goalbox}[1][]{
  colback=gold!8, colframe=gold, boxrule=0.8pt,
  arc=3pt, left=8pt, right=8pt, top=6pt, bottom=6pt,
  breakable, #1
}
\newtcolorbox{notebox}[1][]{
  colback=softgray, colframe=midgray, boxrule=0.6pt,
  arc=3pt, left=8pt, right=8pt, top=6pt, bottom=6pt,
  breakable, #1
}
\newtcolorbox{warnbox}[1][]{
  colback=warnred!6, colframe=warnred, boxrule=0.8pt,
  arc=3pt, left=8pt, right=8pt, top=6pt, bottom=6pt,
  breakable, #1
}
\newtcolorbox{codebox}[1][]{
  colback=codebg, colframe=codebg, coltext=white,
  arc=3pt, left=8pt, right=8pt, top=6pt, bottom=6pt,
  breakable, #1
}

\begin{document}

% ============================================================
\begin{center}
  {\Huge\bfseries\color{gold} خطة إصلاح مطابقة مصطلحات القاموس}\\[4pt]
  {\large\color{darkbg} أن يطابق التعريفُ الكلمةَ في النص بدقة}\\[6pt]
  {\small\color{midgray} منصة يَسِّر الفقهية — \today}
\end{center}

\vspace{4pt}
\hrule
\vspace{10pt}

% ============================================================
\section{المشكلة}

\begin{goalbox}
المطابقة الحالية تبحث عن مصطلح القاموس داخل النص بلا \textbf{حدود كلمات} وبلا \textbf{تسامح صرفي}.
النتيجة: تطابقات خاطئة (مصطلح داخل كلمة أطول)، وفوائت (صيغة في النص لا تطابق مفتاح القاموس)،
ومدير يضطر لتكرار المصطلح بيده بمشكّل وبلا مشكّل.
\end{goalbox}

\subsection*{الدليل المقاس من البيانات}
\begin{itemize}
  \item \textbf{تكرار مصطنع في القاموس} (\texttt{src/data/glossary.ts}): \texttt{مُنْقٍ} و\texttt{منق}،
        \texttt{إداوة} و\texttt{اداوة}، \texttt{عَنَزَة} و\texttt{عنزة}، \texttt{رَجِيع} و\texttt{رجيع}،
        \texttt{البراز} و\texttt{البَراز} — أزواج متطابقة بعد التطبيع، تُثبت أن المطابقة لا تتسامح كفاية.
  \item \textbf{تطابقات خاطئة}: البحث عن \texttt{يد} يطابق داخل «صَعِيدًا»، «سعيد»، «الحديد»، «يدخل».
        و\texttt{غسل} داخل «اغسلني»، والحال نفسه مع أي مصطلح قصير.
  \item \textbf{منطق مكرّر}: المطابقة مكتوبة مرتين — في \texttt{glossaryUtil.tsx} (المولّد القياسي)
        وفي \texttt{QuoteCard.tsx} (نسخة يدوية مطابقة). أي إصلاح في أحدهما لا يسري على الآخر.
\end{itemize}

\begin{notebox}
\textbf{نطاق المطابقة:} متن المسألة (\texttt{QuoteCard}) وشرح الشيخ (\texttt{ExplanationCard})،
عبر \texttt{sortedGlossaryTerms} و\texttt{createGlossaryRenderer} ونسخة \texttt{QuoteCard} المكرّرة.
\end{notebox}

% ============================================================
\section{المرحلة ١ — فرض حدود الكلمات}

\subsection*{المبدأ}
لا يُطابق المصطلح إلا إذا كان \textbf{مستقلاً} (قبله وبعده ليس حرفاً عربياً).

\begin{itemize}
  \item عند كل موضع مرشّح، تحقّق أن الحرف المطبَّع السابق \textbf{ليس} حرفاً عربياً
        (\texttt{[\textbackslash u0600-\textbackslash u06FF]}).
  \item والحرف المطبَّع التالي \textbf{ليس} حرفاً عربياً.
  \item \textbf{استثناء مقصود:} علامات ترقيم/مسافات/نهاية النص تُعامَل كحدود صحيحة.
\end{itemize}

\begin{codebox}
{\ttfamily\small\raggedleft
// المفاهيم (لا الكود النهائي):\\
const isArabic = (c) => /[\textbackslash u0600-\textbackslash u06FF]/.test(c);\\
function atWordBoundary(norm, start, end) \{\\
\hspace*{1em}const beforeOk = start === 0 || !isArabic(norm[start - 1]);\\
\hspace*{1em}const afterOk = end >= norm.length || !isArabic(norm[end]);\\
\hspace*{1em}return beforeOk \&\& afterOk;\\
\}
}
\end{codebox}

\begin{warnbox}
\textbf{الحذر:} فرض الحدود وحده \textbf{يكسر} المصطلحات متعددة الكلمات إذا احتُسبت المسافة بينها.
يجب أن يكون فحص الحدود على \textbf{أول} حرف من المصطلح و\textbf{آخر} حرف منه فقط، مع السماح بالمسافات داخله.
\end{warnbox}

% ============================================================
\section{المرحلة ٢ — التسامح الصرفي (ال التعريف والضمائر)}

\subsection*{المشكلة}
مفتاح القاموس \texttt{الصلاة} لا يطابق نصاً فيه «صلاة» أو «وصلاته» بلا «ال».

\begin{itemize}
  \item \textbf{خيار (أ) — مطابقة اختيارية لـ«ال» التعريف:} عند بناء \texttt{sortedTerms}،
        يُقبل تطابق المصطلح مع أو بدون «ال» البادئة.
  \item \textbf{خيار (ب) — تطبيع المفتاح بتجريده من «ال» مع حفظ الصيغة الأصلية} للنافذة.
  \item \textbf{قاعدة الأمان:} تطبيع جهة المفتاح \textbf{و} جهة النص معاً بنفس الدالة
        (يجري أصلاً في \texttt{normalizeWithIndex})، فلا ينتج انحراف موضع.
  \item \textbf{حدّ التطبيع:} لا قوائم صرفية كاملة (جذر/وزن) — \textbf{مبالغة} وخطأ.
        نكتفي بـ: تشكيل · تطويل · أشكال الألف · ة/ه · ى/ي · «ال» الاختيارية.
\end{itemize}

\begin{notebox}
هذا يحلّ سبب التكرار اليدوي: \texttt{البراز} و\texttt{البَراز} يُدغمان تلقائياً لأن
التشكيل يُطبع و«ال» اختيارية — فيكفي مفتاح واحد.
\end{notebox}

% ============================================================
\section{المرحلة ٣ — توحيد منطق المطابقة}

\begin{warnbox}
\textbf{مشكلة جوهرية:} المنطق مكرّر في موضعين. أي إصلاح في أحدهما يترك الآخر معطوباً.
\end{warnbox}

\begin{itemize}
  \item حصر المطابقة في \textbf{مصدر واحد}: \texttt{createGlossaryRenderer} في \texttt{glossaryUtil.tsx}.
  \item تحويل \texttt{QuoteCard.tsx} لاستدعاء نفس المولّد بدل النسخة اليدوية (سطور 188--246).
  \item التأكد أن كولباكس الأحداث (\texttt{onHover}/\texttt{onClick}/\texttt{onLeave}/\texttt{onMove}) تُمرَّر كما هي.
  \item \textbf{الأثر:} إصلاح واحد يسري على المتن والشرح معاً، ويقلّ الكود المكرّر.
\end{itemize}

% ============================================================
\section{المرحلة ٤ — تنظيف القاموس من التكرار}

\begin{itemize}
  \item بعد إصلاح التسامح، حذف الأزواج المتطابقة بعد التطبيع
        (\texttt{اداوة}، \texttt{عنزة}، \texttt{رجيع}، \texttt{يتخلى}، ...) والإبقاء على نسخة واحدة.
  \item تغيير الفاصل: \textbf{لا يفقد تعريفاً} — تُراجع كل الأزواج يدوياً قبل الحذف عبر سكربت كشف.
  \item \textbf{سكربت مساعد} \texttt{scripts/audit-glossary.mjs}: يطبع أي مفتاحين متساويين بعد التطبيع.
  \item إبقاء \texttt{الطهور}/\texttt{الطهارة} (مصطلحان مختلفان تعريفاً) — التشابه لا يعني التكرار.
\end{itemize}

% ============================================================
\section{المرحلة ٥ — صحة التعريف مع النص}

\begin{itemize}
  \item إضافة \textbf{وضع تحقق} في لوحة الإدارة: عند كتابة مصطلح، يُعرض \emph{عدّ المواضع التي
        سيُبرز فيها في المسائل الحالية} (باستخدام دالة المطابقة نفسها).
  \item \textbf{قيمة مباشرة:} يرى المدير فوراً إن كان المصطلح سيظهر فعلاً أم لا (يكشف الفوائت).
  \item \textbf{تحذير فوري} لو تطابق مصطلح قصير (سطر أو سطران) بكثرة مفرطة (خطر تطابق خاطئ).
  \item \textbf{الملفات:} \texttt{AdminGlossaryTab.tsx} (معاينة المطابقة) + دالة عدّ مشتركة.
\end{itemize}

% ============================================================
\section{الملفات والتحقق}

\textbf{ملفات تُعدَّل:}
\begin{itemize}
  \item \texttt{src/utils/glossaryUtil.tsx} — حدود الكلمات + التسامح الصرفي (المصدر الواحد).
  \item \texttt{src/Components/Content/QuoteCard.tsx} — حذف النسخة المكرّرة واستدعاء المولّد.
  \item \texttt{src/data/glossary.ts} — حذف المفاتيح المكرّرة بعد التطبيع.
  \item \texttt{src/Components/Admin/AdminGlossaryTab.tsx} — معاينة عدد المطابقات.
  \item \texttt{scripts/audit-glossary.mjs} — كشف التكرار (جديد).
\end{itemize}

\vspace{6pt}
\textbf{أوامر التحقق:}
\begin{codebox}
{\ttfamily\small\raggedleft
npm run typecheck\\
npm run lint\\
npm test\\
npm run build
}
\end{codebox}

\subsection*{اختبارات إلزامية جديدة (\texttt{glossaryUtil.test.ts})}
\begin{itemize}
  \item لا تطابق \texttt{يد} داخل «صَعِيدًا»/«سعيد»/«الحديد».
  \item لا تطابق \texttt{ماء} داخل «سماء»/«ماء الورد» (حدّ كلمة).
  \item تطابق \texttt{الصلاة} مع نص فيه «صلاة» و«وصلاته» (تسامح «ال»).
  \item تطابق \texttt{البراز} مع «البَراز» (تسامح تشكيل) بمفتاح واحد.
  \item المصطلح متعدد الكلمات (\texttt{مفهوم المخالفة}) يُطابق بمسافاته الداخلية.
  \item حفاظ المواضع: النص المُبرز يطابق المقطع الأصلي حرفياً.
\end{itemize}

% ============================================================
\section{معايير القبول}

\begin{itemize}
  \item[${\color{okgreen}\faCheckCircle}$] لا تطابق مصطلح داخل كلمة أطول (حدود كلمات صارمة).
  \item[${\color{okgreen}\faCheckCircle}$] «ال» التعريف اختيارية بلا انحراف مواضع.
  \item[${\color{okgreen}\faCheckCircle}$] منطق مطابقة واحد فقط (لا تكرار).
  \item[${\color{okgreen}\faCheckCircle}$] القاموس بلا مفاتيح مكرّرة بعد التطبيع.
  \item[${\color{okgreen}\faCheckCircle}$] لوحة الإدارة تُظهر عدد المطابقات قبل الحفظ.
  \item[${\color{okgreen}\faCheckCircle}$] كل بوابات الجودة خضراء.
\end{itemize}

\end{document}
